\section{The physical pair algebra}\label{nt:algebra}
Fix $N\ge2$, and set
\begin{equation}\label{eq:root-notation}
 d=N-1,\qquad q=e^{\pi\ii/N},\qquad Q=q^2,\qquad
 x=q^d,\qquad \kappa=q^{d^2/2},\qquad \rho=q-q^{-1}.
\end{equation}
There is no restriction to prime or odd $N$. The main strands carry the $N$-dimensional vector space $V_N$; this is what ``complete color'' means at this stage. The operators $H,K,E,F$ encode its quantum-group action. They are concrete matrices, which we specify so that the crossing convention can be checked. One convenient choice of basis is $v_0,\ldots,v_d$ and
\[
 H v_j=(d-2j)v_j,\quad K=q^H,\quad
 E v_j=[j]_q^2v_{j-1},\quad Fv_j=v_{j+1},
 \qquad [j]_q=\frac{q^j-q^{-j}}{q-q^{-1}}.
\]
Missing endpoint vectors are zero. For $k\ge0$, the quantum factorial is
$[k]_q!=\prod_{r=1}^k[r]_q$, with $[0]_q!=1$.
The finite universal crossing, including its Cartan factor, is
\begin{equation}\label{eq:universal-crossing}
 \mathscr R=q^{H\otimes H/2}
 \sum_{k=0}^{d}\frac{q^{k(k-1)/2}\rho^k}{[k]_q!}
 E^k\otimes F^k.
\end{equation}
Every denominator is nonzero. The physical braided crossing is obtained with the fixed exchange and transpose convention; the inverse is the actual matrix inverse. The framing and enhancement are part of the definition, not an unspecified factor appended to an asymptotic formula.

For calculations in the polynomial algebra it is useful to change basis on this same space. We call the new coordinates balanced; they do not define a different invariant. In this basis the operators are
\begin{equation}\label{eq:balanced-module}
 \mathcal E|a\rangle=(Q^{-a}-1)|a-1\rangle,\quad
 \mathcal F|a\rangle=(1-Q^{a+1})|a+1\rangle,\quad
 K|a\rangle=xQ^{-a}|a\rangle.
\end{equation}
The ordinary cup and cap, and the one marked cap, have entries
\begin{equation}\label{eq:cups}
 p_{ab}=Nx^{-b}\delta_{a+b,d},\qquad
 \ell_{ab}=\frac{x^a}{N}\delta_{a+b,d},\qquad
 m_{ab}=\frac{x^b}{N^2}.
\end{equation}
The cup creates two endpoints and the cap joins them. Their two snake identities say that a bend followed by the opposite bend evaluates to a straight strand. The different marked cap implements the reduced invariant and satisfies $mp=1$. The diagonal operator $g=K^{1-N}=(-1)^dK$ is the \emph{pivot}, the correction inserted when an ordinary strand is closed by a trace. For a two-strand operator $A$ commuting with the pair quantum-group action, the original cylindrical evaluation is
\begin{equation}\label{eq:phi-physical}
 \phi(A)=\frac1N\Tr_{V_N\otimes V_N}\bigl[(I\otimes g)A\bigr].
\end{equation}
The factor $1/N$ in this expression is the single cut average. It is not repeated at the auxiliary circles.


In the ordinary highest-weight basis above, put
$p_a=\prod_{j=1}^a(1-Q^j)$. The balanced basis is
$|a\rangle=\rho^a p_a^{-1}v_a$, with
$\mathcal E=\rho E$ and $\mathcal F=\rho F$.
All $p_a$ with $a<N$ are nonzero. This gives an explicit diagonal
change of basis on the same representation. Its factors cancel on
contracted wires. The invariant normalization and crossing comparison
are those of \cite[Theorem~10.1]{Finite}; no asymptotic bound on this
change of basis is used.

\subsection{The full pair algebra}

An operator on two strands must commute with the joint action of the quantum-group generators if it is to represent one of our invariant tangle boxes. The formulas below specify that joint action. The subsequent lemma identifies all such commuting operators; this is the meaning of the term \emph{commutant}.
The balanced pair generators and Casimir are
\begin{equation}\label{eq:pair-casimir}
 \mathcal E_{12}=\mathcal E_1K_2+\mathcal E_2,\quad
 \mathcal F_{12}=\mathcal F_1+K_1^{-1}\mathcal F_2,\quad
 K_{12}=K_1K_2,\quad
 \Omega=\mathcal F_{12}\mathcal E_{12}+qK_{12}+q^{-1}K_{12}^{-1}.
\end{equation}
Let $\Cheb_n(z)=2T_n(z/2)$, where $T_n$ is the usual Chebyshev polynomial. Thus $\Cheb_0=2$, $\Cheb_1=z$ and $\Cheb_{n+1}=z\Cheb_n-\Cheb_{n-1}$. Define
\begin{equation}\label{eq:hermite-algebra}
 P_N(z)=\Cheb_N(z)+2=\prod_{j=0}^{d}(z-c_j),\qquad
 c_j=q^{2j+1}+q^{-2j-1},\qquad
 \cH_N=\CC[z]/(P_N).
\end{equation}
The list of roots includes multiplicity. Since $c_{d-j}=c_j$, all distinct roots are double except the simple root $-2$ at odd $N$.

\begin{lemma}[Full commutant]\label{lem:commutant}
The endomorphisms of $V_N\otimes V_N$ commuting with the three pair generators are exactly $\CC[\Omega]\simeq\cH_N$.
\end{lemma}
\begin{proof}
Let $W_T$ be the span of $|a,T-a\rangle$. The $K_{12}$ eigenvalue one occurs only on $W_d$. The matrices $\mathcal E_{12}:W_{T+1}\to W_T$ for $T<d$, and $\mathcal F_{12}:W_T\to W_{T+1}$ for $T\ge d$, have triangular maximal minors with nonzero diagonals. Consequently the entire pair is generated from $W_d$ by these raising and lowering operators.

On $W_d$, the Casimir is a symmetric tridiagonal matrix with adjacent entries $-(1-Q^{a+1})^2\ne0$. Its first basis vector is cyclic: the $k$th power has last possible coordinate $k$, with coefficient equal to the nonzero product of the first $k$ adjacent entries. An operator commuting with this block is therefore a unique polynomial of degree below $N$ in that block. Subtracting the same polynomial in $\Omega$ from a pair intertwiner gives zero on $W_d$, hence zero everywhere by generation.

Before specialization, use a variable $\zeta$ and the ordinary
highest-weight module of highest weight $d$, with weights $d-2a$,
$E v_a=[a]_\zeta[N-a]_\zeta v_{a-1}$ and $Fv_a=v_{a+1}$.
At the principal root $[N-a]_q=[a]_q$, so this specializes to the
module already defined. The balanced diagonal change is regular there. In the pair, a highest vector of total degree $j$ is obtained by solving the bidiagonal raising equation. If its coefficient at $|a,j-a\rangle$ is $u_a$, the equation at $|a,j-a-1\rangle$ determines $u_{a+1}/u_a$ as the negative ratio of the two nonzero raising entries. Thus for generic $q$ there is one such vector in each degree $j=0,\ldots,d$. Its highest weight is $2d-2j$, its Casimir is $q^{2d-2j+1}+q^{-2d+2j-1}$, and its successive lowering vectors are nonzero through the usual length. The dimensions $2d-2j+1$ sum to $N^2$. Distinct generic Casimirs make these submodules independent.

In degree $T\le d$, precisely the highest vectors $j\le T$ contribute one descendant each. Their eigenvalues give the block characteristic polynomial. At $q=e^{\pi\ii/N}$ it becomes
\[
 \prod_{j=0}^{T}(z-c_j),\qquad c_j=q^{2j+1}+q^{-2j-1}.
\]
All physical matrix entries specialize regularly. The nonzero adjacent entries of the tridiagonal block make a boundary vector cyclic, so this characteristic polynomial is also minimal even when roots coincide. At $T=d$ it is $P_N$. The elementary identity
\[
 \prod_{j=0}^{N-1}(z-q^{2j+1}-q^{-2j-1})=2T_N(z/2)+2
\]
can also be checked by putting $z=w+w^{-1}$: the roots correspond to $w^N=-1$, and the monic degree and leading coefficient agree. The generating property of $W_d$ propagates its vanishing to the full pair.


This proves the polynomial relation and the full commutant assertion.

\end{proof}

At a double root $c$, reduction modulo $(z-c)^2$ leaves
\[
 f(z)=f(c)+f'(c)(z-c)\pmod{(z-c)^2}.
\]
Thus both the value and its first derivative are independent coordinates of the operator. Throughout the paper, retaining the ``full algebra'' or the ``normal data'' means retaining this derivative as well. A polynomial can vanish at every distinct root and still define a nonzero operator in $\cH_N$. Evaluating only at roots would therefore lose part of the invariant.

\subsection{The actual turn and the Newton basis}

We choose a polynomial basis adapted to successive Casimir roots, including their multiplicities. Its normalization is designed so that turning a physical box reverses the basis index exactly. The same basis makes the final closure a coefficient extraction. These are the two facts that will allow us to follow every internal summation index through a large construction.
Put
\begin{equation}\label{eq:newton-basis}
 s_r=\prod_{u=1}^{r}2\sin\frac{\pi u}{N},\quad s_0=1,\qquad
 B_r(z)=\prod_{j=0}^{r-1}(c_j-z),\qquad
 \Bhat_r(z)=\frac{B_r(z)}{s_r^2},\quad 0\le r\le d.
\end{equation}
The quarter-turn is the literal contraction
\[
 \Turn(A)=(I^{\otimes2}\otimes\ell)(I\otimes A\otimes I)
 (p\otimes I^{\otimes2}).
\]
Its entries in the recorded all-primal chart are
\[
 \Turn(A)_{(a,b);(c,e)}=x^{a-e}A_{(b,d-e);(d-a,c)}.
\]
The cups and caps are those of \eqref{eq:cups}. Direct substitution
in their entries gives
\[
 \mathcal E_{12}p=\mathcal F_{12}p=0,\quad K_{12}p=p,
 \qquad
 \ell\mathcal E_{12}=\ell\mathcal F_{12}=0,\quad \ell K_{12}=\ell.
\]
Together with the two snake identities, these equations show that
contracting a pair intertwiner by the displayed cup and cap again
commutes with the pair generators: move each coproduct summand through
the intertwiner and cancel the cup and cap terms. Thus the turned
operator belongs to the polynomial commutant proved above. The turn
used here has no additional scalar rescaling.

\begin{proposition}[Newton turn and closure]\label{prop:newton-turn}
In the full algebra,
\begin{equation}\label{eq:newton-turn}
 \Turn\Bhat_r=\Bhat_{d-r},\qquad
 \phi\left(\sum_{r=0}^{d}\alpha_r\Bhat_r\right)=\alpha_d,
 \qquad
 \phi(\Bhat_i\Bhat_j)=\delta_{i+j,d}.
\end{equation}
\end{proposition}
\begin{proof}
Now $B_r(\Omega)$ annihilates every degree block $W_T$ with $T<r$, with multiplicities, and has first-color bandwidth at most $r$. In the turned matrix its first-color displacement is $d-(a+b)$, so $\Turn B_r$ annihilates every $W_T$ with $T<d-r$. Its representing polynomial is therefore divisible by $B_{d-r}$.

For the reverse degree bound, its first middle-block column is
\[
 (\Turn B_r)_{a0}=x^{a-d}(B_r(\Omega))_{(d-a,0);(d-a,0)},
\]
which vanishes for $a>d-r$. A polynomial of degree $k<N$ in the cyclic tridiagonal middle block has a nonzero $k$th coordinate on the first basis vector, from its unique outermost path. Thus $\deg\Turn B_r\le d-r$, proving $\Turn B_r=C_rB_{d-r}$.

Use the physical entry from $(d-r,0)$ to $(0,d-r)$. On the turned side it is the unique $r$-step path in the high-degree block, whose factors are
\[
 Q^{j-r}(1-Q^{r-j})(1-Q^{j+1}),\qquad0\le j<r.
\]
Including the leading coefficient $(-1)^r$ gives $s_r^2$. The corresponding path for $B_{d-r}$ in the low block gives $s_{d-r}^2$. All turn monomials in this entry equal one. Hence
\[
 \Turn B_r=\frac{s_r^2}{s_{d-r}^2}B_{d-r},
 \qquad \Turn\Bhat_r=\Bhat_{d-r}.
\]
Finally $s_rs_{d-r}=N$ and the closure calculation immediately below give
$\phi(\Bhat_i\Bhat_j)=\delta_{i+j,d}$: use degree below $d$, leading coefficient at $d$, and full divisibility by $P_N$ above it. This proves Proposition~\ref{prop:newton-turn} by physical support and exact paths.

\end{proof}
The turn is linear and involutive in this basis, but it is not multiplicative.


We next identify the physical trace with its polynomial functional.
This identification is needed in addition to the spectral relation.
Use a generic variable $\zeta$ near $q$, keeping the polynomial $f$
fixed, with balanced module entries
\[
 K_a=\zeta^{d-2a},\quad
 \mathcal F|a\rangle=(1-\zeta^{2a+2})|a+1\rangle,\quad
 \mathcal E|a\rangle=(\zeta^{-N}-\zeta^{N-2a})|a-1\rangle.
\]
They specialize to \eqref{eq:balanced-module}. Let
$C(\zeta)=\operatorname{Tr}_2((I\otimes K)f(\Omega_\zeta))$.
Tracing the $\mathcal E$ coproduct commutator gives
$[C(\zeta),\mathcal E]=0$. For the lowering generator, conservation
of $K_1K_2$ gives, with $A=f(\Omega_\zeta)$,
\[
 A K_1^{-1}=K_1^{-1}K_2^{-1} A K_2,
\]
so cyclicity of the second partial trace yields
\[
 \operatorname{Tr}_2 K_2[A,K_1^{-1}\mathcal F_2]
 =K_1^{-1}\operatorname{Tr}_2
       (A K_2\mathcal F_2-K_2\mathcal F_2 A)=0.
\]
The remaining coproduct term gives $[C(\zeta),\mathcal F]=0$.
The distinct $K$ weights and nonzero adjacent lowering entries imply
that $C(\zeta)$ is a scalar, which we denote by $c_f(\zeta)$.
Taking one more $K$ trace and using the generic highest-weight
decomposition proved above gives
\begin{equation}\label{nt:generic-trace}
 c_f(\zeta)=
 \frac{\displaystyle\sum_{j=0}^{d}[2N-1-2j]_\zeta
 f(\zeta^{2N-1-2j}+\zeta^{-2N+1+2j})}{[N]_\zeta}.
\end{equation}
The finite matrices show that this quotient has a removable value at
$q$. Differentiate numerator and denominator with respect to
$\log\zeta$, keeping $f$ fixed. The denominator derivative is
$-2N/\rho$. For a pair of indices $m,2N-m$ with common limiting
node $c=q^m+q^{-m}$, the two dimension weights cancel at $q$;
their derivative sum is $2Nc/\rho$. The sum of dimension weight
times node derivative is $2N(c^2-4)/\rho$. This pair therefore
contributes $(4-c^2)f'(c)-cf(c)$ to the removable quotient.
When $N$ is odd the self-paired index $m=N$ contributes $f(-2)$.
At the fixed root the original pivot is $g=(-1)^dK$, and the
normalizing factor $1/N$ in \eqref{eq:phi-physical} cancels the
trace of the scalar operator on the remaining $N$-dimensional head.
Consequently the actual closure is

\begin{equation}\label{eq:hermite-trace}
 \phi(f)=(-1)^d\sum_{c\ \mathrm{double}}
 \bigl((4-c^2)f'(c)-cf(c)\bigr)
 +\mathbf1_{N\ \mathrm{odd}}f(-2).
\end{equation}
Equivalently,
\begin{equation}\label{eq:leading-trace}
 \phi(f)=(-1)^dN^2[z^d](f\bmod P_N).
\end{equation}
For example, at a double root $c$, writing $Q_c=P_N/(z-c)^2$ gives
$Q_c(c)=N^2/(4-c^2)$ and $Q_c'(c)/Q_c(c)=c/(4-c^2)$ from the Chebyshev differential equation. The residue of $f/P_N$ is exactly the corresponding term of \eqref{eq:hermite-trace} after multiplication by $(-1)^dN^2$. Summing residues proves \eqref{eq:leading-trace}; the simple odd root has $P_N'(-2)=N^2$. The sole shortest tridiagonal path from the first to the last coordinate has product $(-1)^dN^2$, which also identifies this scalar with the original endpoint matrix element.

\subsection{Positive multiplication with every normal retained}

Suppose two stacked boxes contribute basis indices $i$ and $j$. Their product is a linear combination of output indices $r$. The following coefficient formula gives both the allowed triples $(i,j,r)$ and their weights. Calling this local operation a \emph{gate} is only shorthand for this one multiplication; the geometric operation remains the vertical stack.
\begin{theorem}[Multiplication law]\label{thm:positive-product}
If $i+j<N$, then
\begin{equation}\label{eq:positive-product}
 \Bhat_i\Bhat_j=\sum_{r=\max(i,j)}^{i+j}\mu_{ijr}\Bhat_r,
 \qquad
 \mu_{ijr}=\frac{s_{i+j}s_r}
 {s_i s_j s_{r-i}s_{r-j}s_{i+j-r}}>0.
\end{equation}
If $i+j\ge N$, the product is zero in the full algebra. In particular, a product gate has precisely the integer support
\begin{equation}\label{eq:gate-support}
 i+j\le d,\qquad \max(i,j)\le r\le i+j.
\end{equation}
\end{theorem}
\begin{proof}
At a generic parameter $\zeta$, put $t=-(\zeta-\zeta^{-1})^2$. The unnormalized coefficients are
\[
 B_iB_j=\sum_r
 t^{i+j-r}\frac{[i]![j]![i+j]!}
 {[r]![r-i]![r-j]![i+j-r]!}\,B_r.
\]
Multiplication by $c_i-z$ and the identity
$(c_i-z)B_r=B_{r+1}+(c_i-c_r)B_r$, with
$c_i-c_r=t[r+i+1][r-i]$, reduce induction to
\[
 [i+1][i+j+1]=[r][r-j]+[r+i+1][i+j+1-r].
\]
This follows by expanding the four powers of $\zeta$. Endpoint coefficients follow from the same zero factors. At the principal root, all factorial indices in the case $i+j<N$ are below $N$. Since $[u]!=s_u/(2\sin(\pi/N))^u$, normalization by $s_i^2s_j^2/s_r^2$ gives \eqref{eq:positive-product}.

For $i+j\ge N$, use the exact polynomial identity $B_N/B_j=B_{N-j}$, a consequence of $c_{d-u}=c_u$. Then $B_iB_j$ is divisible by $B_N=(-1)^NP_N$ with multiplicities. This proves the zero assertion without dropping the first normal directions.
\end{proof}

\section{Colored circles and retained half twists}\label{nt:colors}

A color $j$ on an auxiliary circle means that the circle carries a $j$-dimensional representation. The main strands still carry $V_N$. We must compute the resulting operator on those two main strands, including the trace around the auxiliary circle, before replacing that circle by twists. The answer is a Chebyshev polynomial in $\Omega$, with the sign fixed by the pivot.
For $1\le j\le N$, the auxiliary module is the ordinary $j$-dimensional highest-weight module. Its loop is the actual block-braiding word $[2,1,1,2]$, with the intermediate strand dimensions retained, followed by its ordinary pivot $K_j^{1-N}$. There is no auxiliary $1/j$ or $1/N$.

\begin{proposition}[Circle identity]\label{prop:circle}
The resulting two-head operator is
\begin{equation}\label{eq:circle}
 L_j=(-1)^{j-1}U_{j-1}(\Omega/2).
\end{equation}
Here $U_0=1$, $U_1(z/2)=z$, and $U_{n+1}(z/2)=zU_n(z/2)-U_{n-1}(z/2)$.
\end{proposition}
\begin{proof}
First use a generic parameter and the ordinary $K_j$ pivot on the
auxiliary module. Evaluate the resulting closed monodromy on a target
highest vector of weight $n$. In one direction the positive raising terms annihilate that vector; in the return direction a positive lowering term cannot return to its highest coefficient. The two Cartan factors and this temporary $K_j$ pivot give
\[
 \sum_{a=0}^{j-1}q^{(n+1)(j-1-2a)}
 =U_{j-1}\bigl((q^{n+1}+q^{-n-1})/2\bigr).
\]
The generic pair highest vectors have distinct Casimirs and descendants of total dimension $N^2$, so this proves the full matrix identity. All finite crossing denominators have indices below $N$, so the
$K_j$ identity specializes regularly. Only after specialization replace
that pivot by the actual $K_j^{1-N}$: at the root,
$K_j^{1-N}=(-1)^{j-1}K_j$. This proves \eqref{eq:circle}, including its nilpotent coefficients.
\end{proof}

In particular $L_N=(-1)^dP_N'(\Omega)/N$. The trace identity
\begin{equation}\label{eq:circle-trace}
 \phi(L_Nf)=N\sum_{j=0}^{d}f(c_j)
\end{equation}
counts the repeated roots. It follows first by Lagrange interpolation for a polynomial with distinct roots and then by polynomial continuation in its coefficients. At even levels $L_N^2=0$; at odd levels $L_N^2=N^2E_{-2}$, with the simple-node projector $E_{-2}$. Thus complete color alone cannot be inverted to reconstruct a filling.

Let $M=\lfloor d/2\rfloor$ and $h_j=\Turn L_j$. The complete leaf has the exact positive expansion
\begin{equation}\label{eq:complete-leaf}
 h_N=\sum_{v=0}^{M}a_{N,v}\Bhat_v,
 \qquad a_{N,v}=\frac{N^2s_{2v}}{s_v^4}>0.
\end{equation}
For $0\le h\le M$, with $h_0=0$, the full color decomposition is
\begin{equation}\label{eq:core-alias}
 h_{N-h}=\sum_{v=h}^{M}
 \frac{N^2s_{2v}}{s_v^2s_{v-h}s_{v+h}}\Bhat_v+(-1)^Nh_h.
\end{equation}
For $1\le h\le\lfloor N/2\rfloor$,
\begin{equation}\label{eq:low-alias}
 h_h=\sum_{r=0}^{h-1}(-1)^{h-1+r}
 \frac{s_r^2s_{h+r}}{s_{2r+1}s_{h-r-1}}\Bhat_{d-r}.
\end{equation}
These identities separate two different ranges of Newton indices. The first term in \eqref{eq:core-alias}, called the \emph{core}, has positive coefficients before a half twist and occupies the lower indices. The second term, called the \emph{alias}, is an exact signed contribution from the complementary color; formula~\eqref{eq:low-alias} puts it in the higher Newton indices. An alias is therefore part of the finite invariant, not an error introduced by approximation. The following polynomial derivation proves the formulas with their full derivative data. Controlling the aliased contributions after arbitrary compositions requires further estimates.

\subsection{Proof of the color expansions}
We derive the coefficient formulas rather than infer them from values at the distinct Casimir roots. That distinction is necessary because the quotient algebra also contains first derivatives at double roots. The generating-function calculation first gives a polynomial identity at a generic parameter; regular specialization then retains all those derivatives automatically.
Here $[k]_q!$ is the quantum factorial and $\genfrac{[}{]}{0pt}{}{a}{b}_q=[a]_q!/([b]_q![a-b]_q!)$ at a generic parameter. The generic Chebyshev--Newton identity is
\begin{equation}\label{eq:A-Cheb-newton}
 U_{j-1}(z/2)=\sum_{r=0}^{j-1}(-1)^r
 \genfrac{[}{]}{0pt}{}{j+r}{2r+1}_q B_r(z).
\end{equation}
To prove it, successively divide
$1-zt+t^2=(1-c_rt+t^2)+t(c_r-z)$. The formal identity is
\[
 \frac1{1-zt+t^2}
 =\sum_{r\ge0}\frac{(-t)^rB_r(z)}{\prod_{i=0}^{r}(1-c_it+t^2)}.
\]
After $r$ steps the remainder is divisible by $t^{r+1}$. The denominator factors as a consecutive $q$-string, and the Gaussian Pascal recurrence gives its reciprocal coefficient. Taking the coefficient of $t^{j-1}$ gives \eqref{eq:A-Cheb-newton}. For $j\le\lfloor N/2\rfloor$, every factorial index is below $N$; specializing and turning proves \eqref{eq:low-alias}.

For higher colors use the polynomial identity
\[
 U_{N-h-1}=\tfrac12(\Cheb_hU_{N-1}-\Cheb_NU_{h-1}).
\]
Since $\Cheb_N=-2$ in the full algebra,
\[
 L_{N-h}=(-1)^h\tfrac12\Cheb_hL_N+(-1)^NL_h.
\]
The $v$th turned coefficient of the first term is
$(-1)^hN\sum_{k=0}^{d}\cos(h(2k+1)\pi/N)\Bhat_v(c_k)$.
This finite sum can be evaluated without a zero denominator. Put $z=w+w^{-1}$. The factorization of $B_v$ is the consecutive string
\[
 B_v(w+w^{-1})=(-1)^vw^{-v}
 \prod_{j=0}^{2v-1}\left(w-q^{-(2v-1)+2j}\right).
\]
For $0\le v,h\le M$, all its Fourier exponents after multiplication by $w^{\pm h}$ have absolute value below $N$. Summing at $w=q^{2k+1}$ extracts only the constant coefficient. The finite Gaussian product gives that coefficient as zero for $v<h$, and as
$(-1)^h s_{2v}/(s_{v-h}s_{v+h})$ for $v\ge h$. Dividing by $s_v^2$ and multiplying the two factors of $N$ proves \eqref{eq:core-alias}, including $h=0$ and both parities.


\subsection{A retained half twist}

The preceding formulas describe a circle with no retained crossing. We next insert one positive or negative crossing in its original disk framing. This changes both the coefficients and their phases. The sign $\eps$ refers to this one physical crossing in its disk framing.
Write $R$ for the physical crossing of
\cite[Proposition 5.1]{Finite}, in the balanced basis. Its scalar is
$\kappa=q^{d^2/2}$ with the lift in \eqref{eq:root-notation}.
A retained half twist acts on a turned leaf by
$\Turn_\eps(f)=\Turn(R^\eps\Turn f)$, for $\eps=\pm1$. Its exact transport is
\begin{equation}\label{eq:half-twist}
 \Turn_\eps\Bhat_v=\kappa^\eps(-1)^v
 \sum_{r=0}^{v}q^{\eps v(r+1)}
 \frac{s_v}{s_rs_{v-r}}\Bhat_r.
\end{equation}
The triangular formula retains the full Cartan phase and all normals. Its derivation uses the finite continuant expansion and root factorial reflection; the two signs are coefficient conjugates in the real Newton basis. A half-twisted alias may occupy every Newton coordinate, so the pure alias support cannot be reused after \eqref{eq:half-twist}.


Here is a proof identifying the entire crossing polynomial. In this
paragraph use a generic parameter $\zeta$ and put
\[
 \sigma_r=(\zeta^r-\zeta^{-r})/\ii,\quad
 s_r=\prod_{a=1}^r\sigma_a,\quad
 c_r=\zeta^{2r+1}+\zeta^{-2r-1}.
\]
The matrix of multiplication by $z$ in the normalized Newton basis
is lower bidiagonal, with $Z_{rr}=c_r$ and
$Z_{r+1,r}=-\sigma_{r+1}^2$. Define the triangular matrix
\[
 A_{ij}=(-1)^i\zeta^{i(j+1)}\frac{s_i}{s_js_{i-j}},
 \qquad i\ge j,
\]
and set its other entries to zero. It commutes with $Z$.
Indeed, after division by $A_{ij}$ for $i>j$, the only entry to
check is
\[
 c_i-c_j+\sigma_{i-j}
       (\zeta^{-j-1}\sigma_i+\zeta^i\sigma_{j+1})=0.
\]
The parenthesis is $\sigma_{i+j+1}$ and
$c_i-c_j=-\sigma_{i+j+1}\sigma_{i-j}$. The diagonal and boundary
entries follow from the same identities with zero factors retained.
The generic diagonal nodes are distinct, so $A$ is the polynomial
in $Z$ whose values are $\nu_j=(-1)^j\zeta^{j(j+1)}$.
Its first column gives that polynomial explicitly:
\[
 r_\zeta(z)=\sum_{r=0}^d(-\zeta)^r\Bhat_r(z).
\]
Every coefficient is regular near $q$, since its factorial indices
are below $N$. At a paired root $c_j=c_{d-j}$ the divided difference
of the generic values has limit
\[
 \lim_{\zeta\to q}
 \frac{\nu_j-\nu_{d-j}}{c_j-c_{d-j}}
 =-\nu_j\frac{d-2j}{2(q^{2j+1}-q^{-2j-1})}.
\]
For example, differentiating in $\log\zeta$ gives numerator
$-N(d-2j)\nu_j$ and denominator
$2N(q^{2j+1}-q^{-2j-1})$. These are exactly the first derivatives
of the actual crossing polynomial in \cite[Proposition 5.1]{Finite},
in addition to its values. The odd simple root requires only its
value. Uniqueness of Hermite interpolation therefore identifies
\begin{equation}\label{nt:crossing-newton}
 R=\kappa\sum_{r=0}^d(-q)^r\Bhat_r(\Omega),\qquad
 (R(Z))_{ij}=\kappa(-1)^i q^{i(j+1)}
                           \frac{s_i}{s_js_{i-j}}\quad(i\ge j).
\end{equation}

Reversal of both basis indices gives the displayed half-twist formula.
Explicitly $s_a s_{d-a}=N$, and
\[
 (d-r)(d-v+1)-v(r+1)=N(d-r-v).
\]
Thus the reflected scalar is $\kappa(-1)^v q^{v(r+1)}$, with
factorial ratio $s_v/(s_rs_{v-r})$. Replacing $q$ by $q^{-1}$
conjugates the coefficients of the polynomial while keeping the real
Newton basis fixed. The Hermite values and logarithmic derivatives
then give the inverse crossing, proving the formula for both signs.

